\subsection{Literacy, Training, and Where Each One Sits}
\label{sec:8.3.5}
A safeguard nobody understands is a safeguard nobody can invoke, which is what puts AI literacy — among developers, policymakers, and the public alike — on the list of things the rest of this chapter's proposals quietly assume. Section~\ref{sec:5.4.1} names cross-training among developers, ethicists, and legal scholars as a lever on a system's immediate environment. Both depend on a question the phrases hide, which is where the teaching sits.

\runin{The public tier is the one nobody owns}
Developers and students are the easy case, and the mechanisms are familiar: AI concepts and their ethical stakes introduced before higher education, interdisciplinary degree programs, specialized training for working engineers that goes past a compliance checklist. All of it reaches people already inside the pipeline or heading for it, and it works about as well as curriculum reform generally works.

The public is hard because no institution has the job. Elements of AI, a free course built jointly by the University of Helsinki and the Finnish company Reaktor, is the clearest instance of somebody treating public understanding as an obligation rather than a market: Finland funded its translation into every EU member-state language during its EU presidency, specifically to carry it past its own citizens \autocite{helsinki2020elements}. That took a national government deciding a course was infrastructure. Nothing in the ordinary funding of AI education produces that decision, which is why the public tier stays thin while the developer tier fills in.

Education for communities historically locked out of the field is part of this argument and not a separate one about diversity. A next generation of developers drawn from a narrow slice of the population reproduces exactly the narrow set of assumptions section~\ref{sec:6.1.1} identifies as a source of AI bias, which makes access a technical requirement on the systems as much as a fairness claim about who builds them.

\runin{Elective or embedded}
For ethics training the real choice is between a course a student can decline and a requirement they cannot. Embedded EthiCS, developed at Harvard by the computer scientist Barbara Grosz and the philosopher Alison Simmons, pairs philosophy graduate students with computer-science faculty to write ethics modules directly into standard CS courses \autocite{grosz2019embedded}. The placement is the design. A stand-alone ethics class is taken by the students who elect it, and its content stays separable from the technical work in a student's mind because it was separable in the timetable. A module inside the algorithms course is not separable from the algorithms. Grosz's own course works the same way at the level of content, built on cases with consequences — predictive policing, public-health interventions, wildlife conservation — instead of abstract dilemmas \autocite{harvardgazette2019ethics,harvardcs108schedule}.

Certification does the equivalent work for people whose formal education is finished. The IEEE's CertifAIEd program builds a professional credential around accountability, transparency, privacy, and bias, open to developers and equally to the policymakers, underwriters, and human-resources staff whose work touches an AI system without building one \autocite{ieeesa2020certifaied}. A credential is a weak instrument: it certifies that somebody sat the assessment, and no more. It is also the only instrument here that reaches anyone who has already left a university, which is most of the people making decisions about these systems today.
